% TOP_PROBLEM: {"problemId": "problem.top500-omission-w045049-stationary-integral-varifold-singular-set", "canonicalTitle": "Measure-zero singular set of stationary integral varifolds", "releaseRank": 414, "primaryDomain": "analysis_pde", "primaryDomainLabel": "Analysis & PDE", "displayStatus": "open_with_solved_subcases", "releaseStatus": "open_with_solved_subcases"}
% !TeX program = lualatex
% Definition and sources only; this file is not an LLM solution attempt.
\documentclass[11pt]{article}
\usepackage[margin=25mm]{geometry}
\usepackage{fontspec}
\setmainfont{Latin Modern Roman}
\usepackage{amsmath,amssymb,mathtools}
\usepackage{unicode-math}
\setmathfont{Latin Modern Math}
\usepackage{xurl,hyperref}
\hypersetup{colorlinks=true,urlcolor=blue}
\setcounter{secnumdepth}{0}
\setlength{\parindent}{0pt}
\setlength{\parskip}{5pt}
\setlength{\emergencystretch}{3em}
\begin{document}
\section*{\texorpdfstring{Measure-zero singular set of stationary integral varifolds}{Measure-zero singular set of stationary integral varifolds}}\label{414}
\subsection{Definitions and mathematical statement}
An integral $m$-varifold in $U\subset{\mathbb{R}}^{m+n}$ is a Radon measure on position--tangent-plane pairs represented by a countably $m$-rectifiable set with locally integrable positive integer multiplicity. It is stationary when its first variation vanishes:
\[
 \int\operatorname{div}_P X(x){\,\mathrm{d}} V(x,P)=0
 \quad\text{for every }X\in C_c^1(U;{\mathbb{R}}^{m+n}).
\]
A regular point of $\operatorname{spt}V\cap U$ has a neighborhood on which $V$ is an integer multiple of a smooth embedded minimal $m$-submanifold. Let $\operatorname{Sing}(V)$ be the remaining support points. The conjecture is $\mathcal H^m(\operatorname{Sing}(V))=0$ for all positive $m,n$. Stationarity alone is assumed: area minimization, stability, and multiplicity one are not extra hypotheses. Rectifiability almost everywhere is weaker than having these open smooth neighborhoods.
\par\smallskip\textit{Formulation and concept references:} \href{https://arxiv.org/html/2503.00651}{[S1]}.

\subsection{Short English statement}
Does every stationary integral varifold have a singular set of zero measure in its own dimension, without any stability or minimizing assumption?
\subsection{Sources}
\begin{itemize}
\item[S1]Brena, Decio and De Lellis, Remarks and Conjectures on Stationary Varifolds (2025).
\url{https://arxiv.org/html/2503.00651}.
\textit{Formulation source.}
\end{itemize}
\end{document}
